% \iffalse meta-comment
%
% hit-thesis-meta.dtx 是学位论文的元信息字段
%
% \fi
%
% \section{学位论文的元信息字段}
%
% \changes{v3.2a}{0000/00/00}{元信息字段的声明单独成文件。它说的是「有哪些字段」，
%   与 \file{src/config/} 说的「字段取什么值」是一对，都是数据；
%   字段怎么存、旧名怎么兼容那套机制在 \file{src/utils/hit-meta.dtx}}
% 学位论文用到的元信息字段：字段名、旧名映射、\cs{hitsetup} 各族的路由，
% 以及摘要环境。都是清单，不含机制。
%
% ^^A ======================================================================
% ^^A Module thesis-meta: 元信息字段（\hitsetup 接受的 key）（hit-thesis）
% ^^A ======================================================================
%    \begin{macrocode}
%<*thesis-meta>
% \changes{v3.2a}{0000/00/00}{meta 模块抽取后，etoolbox 需提前加载以使用 \cs{AfterPreamble}}
\RequirePackage{etoolbox}
% \changes{v3.2a}{0000/00/00}{\cs{hitsetup} 的底层机制挪到 \file{src/utils/hit-meta.dtx}}

\hit@define@term{secret-level} %密级
\hit@define@term{classified-index-national}  %国内图书分类号
\hit@define@term{classified-index-international}  %国际图书分类号

\hit@define@title{zh}
\hit@define@term{discipline-zh} %中文学科
\hit@define@term{author-zh} %中文作者
\hit@define@term{supervisor-zh} %中文导师
\hit@define@term{associate-supervisor-zh} %中文副导师
\hit@define@term{co-supervisor-zh}%中文联合导师
\hit@define@term{affiliation-zh}%中文院系
\hit@define@term{specialty-zh}%中文专业
%    \end{macrocode}
%
% \changes{v3.1b}{2022/06/06}{增加深圳封面最下方的时间信息}
% \changes{v3.1d}{2025/03/03}{根据 \issue[226]，将 \texttt{szshortcdate} 改为 \texttt{firstpagecdate}，只要赋值就会显示}
% 根据 \issue[95] 增加一个变量 \texttt{szshortcdate} 来存储答辩时间。
% 根据 \issue[226]，将 \texttt{szshortcdate} 改为 \texttt{firstpagecdate}，只要赋值就会显示。
%    \begin{macrocode}
\hit@define@term{cover-date-zh}
\hit@define@term{date-zh}
\hit@define@term{report-number}
\hit@define@term{station-zh}
%    \end{macrocode}
%
% 添加博后封皮选项
% \changes{v3.0.12}{2020/11/02}{添加博后封皮选项}
%    \begin{macrocode}
\hit@define@term{work-finish-date-zh}
\hit@define@term{report-submit-date-zh}
\hit@define@term{research-start-date-zh}
\hit@define@term{research-end-date-zh}
%    \end{macrocode}
%
% 添加博后封皮的title选项
% \changes{v3.0.0}{2020/05/21}{添加博后封皮的title选项}
%    \begin{macrocode}

\hit@define@term{student-id}%
\hit@define@term{student-type-zh}%



\hit@define@title{en}
\hit@define@term{discipline-en} %英文学科
\hit@define@term{author-en} %英文作者
\hit@define@term{supervisor-en} %英文导师
\hit@define@term{associate-supervisor-en} %英文副导师
\hit@define@term{co-supervisor-en} %英文联合导师
\hit@define@term{affiliation-en}
\hit@define@term{specialty-en}
\hit@define@term{date-en}
\hit@define@term{student-type-en}
\cs_new_protected:Npn \hit@collect@abstract@zh #1 { \cs_gset:Npn \hit@abstract@zh {#1} }
\cs_new_protected:Npn \hit@collect@abstract@en #1 { \cs_gset:Npn \hit@abstract@en {#1} }
\NewDocumentEnvironment { abstract-zh } { } { \Collect@Body \hit@collect@abstract@zh } { }
\NewDocumentEnvironment { abstract-en } { } { \Collect@Body \hit@collect@abstract@en } { }
\NewDocumentEnvironment { cabstract } { }
  {
    \hit@warn@renamed@environment { cabstract } { abstract-zh }
    \Collect@Body \hit@collect@abstract@zh
  }
  { }
\NewDocumentEnvironment { eabstract } { }
  {
    \hit@warn@renamed@environment { eabstract } { abstract-en }
    \Collect@Body \hit@collect@abstract@en
  }
  { }
%    \end{macrocode}
%
% \changes{v3.2a}{0000/00/00}{关键词的逗号拆分挪到 \file{src/utils/hit-meta.dtx}}
% 关键词的拆分机制见 \file{src/utils/hit-meta.dtx}。
%    \begin{macrocode}
\hit@parse@keywords{keywords-zh}{keywords@separator@zh}
\hit@parse@keywords{keywords-en}{keywords@separator@en}
\keys_define:nn { hit / class }
  {
    design .code:n =
      {
        \str_if_eq:nnTF {#1} { true }
          { \bool_gset_true:N  \g__hit_design_bool }
          { \bool_gset_false:N \g__hit_design_bool }
      } ,
    design .default:n = { true } ,
  }
%    \end{macrocode}
%
% \changes{v3.2a}{0000/00/00}{\cs{hitsetup} 改用 l3keys 做入口，加一层按职责
%   划分的子族路由。子族只做转发，键由 \texttt{hit/class} 族处理，行为
%   与改动前一致；后续再逐族把键迁过来}
% \cs{hitsetup} 的入口换成 l3keys。根族只做路由，把 \option{info}、\option{style}
% 等分派到对应子族；子族一律把键原样转回 \texttt{hit/class}
% 族，所以两种写法结果相同：
% \begin{latex}
% \hitsetup{title-zh={...}}              % 扁平写法，一直支持
% \hitsetup{info={title-zh={...}}}       % 新的分族写法
% \end{latex}
% 根族的 \option{unknown} 兜底负责扁平写法，子族的兜底负责分族写法。等某个键真正
% 迁到 l3keys 之后，它就不再走兜底，两条路自然收敛。兜底本身、三态取值与布尔键的
% 声明机制都在 \file{src/utils/hit-meta.dtx}。
%
% \changes{v3.2a}{0000/00/00}{元信息字段正式注册进 \texttt{hit/info} 族，
%   不再依赖子族兜底转发}
% 元信息字段名单。这批由 \cs{hit@define@term} 与 \cs{hit@parse@keywords} 批量生成，
% 每个字段同时有一个同名命令（\cs{title-zh} 等），键的处理器直接调用那个命令，
% 所以三种写法结果一致：\cs{hitsetup} 的扁平写法、分族写法、以及命令写法。
%    \begin{macrocode}
\hit@declare@info@field
  {
    secret-level , classified-index-national , classified-index-international , title-zh ,
    discipline-zh , author-zh , supervisor-zh , associate-supervisor-zh , co-supervisor-zh ,
    affiliation-zh , specialty-zh , cover-date-zh , date-zh , report-number , station-zh ,
    work-finish-date-zh , report-submit-date-zh , research-start-date-zh , research-end-date-zh
    , student-id , student-type-zh , title-en , discipline-en , author-en , supervisor-en ,
    associate-supervisor-en , co-supervisor-en , affiliation-en , specialty-en , date-en ,
    student-type-en , keywords-zh , keywords-en
  }
%    \end{macrocode}
%
% \changes{v3.2a}{0000/00/00}{\file{cfg} 里的可本地化字符串做成 \texttt{hit/const}
%   族的键，用户在导言区即可改写，不必去动生成物}
% 可本地化字符串的名单。默认值在类文件末尾给出，\cs{hitsetup} 在导言区执行，
% 所以用户设的值覆盖默认值。键名怎么由宏名得来见 \file{src/utils/hit-meta.dtx}。
%    \begin{macrocode}
\hit@declare@constant@as
  {
    hithesis , hit , pozhehao , PGR , UGR ,
    equationname , listfigureename , listtableename , listequationename ,
    listequationname , cabstractcname , cabstractename , eabstractcname , eabstractename
  }
\hit@declare@constant@ctex
  {
    chapter-name = chapter/name , appendixname , contentsname , listfigurename ,
    listtablename , figurename , tablename , bibname , indexname
  }
\hit@declare@constant
  {
    bachelor-declaration-section , bachelor-authorization-section ,
    bachelor-declaration-heading-line-one , bachelor-declaration-heading-line-two ,
    bachelor-declaration-heading-short ,
    algorithm-name-zh , algorithm-name-en ,
    campus-full ,
    acknowledgement-heading-zh , acknowledgement-heading-en , authorization-heading-zh ,
    authorization-heading-en , authorization-text , bachelor-authorization-text ,
    authorization-usage-heading , author-signature-field-label ,
    bachelor-affiliation-field-label , bachelor-date-field-label ,
    bachelor-check-date-field-label , bachelor-student-field-label ,
    bachelor-document-type , bachelor-degree-level ,
    bachelor-teacher-comment-field-label , bachelor-teacher-signature-field-label ,
    bibliography-heading-en , bibliography-heading-zh , affiliation-field-label-zh , associate-supervisor-field-label-zh ,
    author-field-label-zh , brace-left-zh , brace-right-zh , co-supervisor-field-label-zh ,
    date-field-label-zh , degree-of-discipline-format-zh , degree-field-label-zh ,
    keywords-separator-zh , keywords-field-label-zh , conclusion-heading-zh ,
    keywords-field-label-en , toc-page-heading , toc-table-heading ,
    toc-figure-heading , nomenclature-heading ,
    conclusion-heading-en , corresponding-address-heading-zh , corresponding-address-heading-en , school-zh , school-field-label-zh
    , specialty-field-label-zh , supervisor-field-label-zh , document-type ,
    thesis-title-prefix , degree-level-zh , date-blank , declaration-heading ,
    bachelor-declaration-heading , declaration-text , denotation-heading-zh ,
    denotation-heading-en , eqdenote-label-zh , eqdenote-label-en , eqdenote-dash ,
    doctor-publication-heading-zh , doctor-publication-heading-en , affiliation-field-label-en ,
    associate-supervisor-field-label-en , author-field-label-en , brace-left-en , brace-right-en
    , co-supervisor-field-label-en , date-field-label-en , degree-of-discipline-format-en ,
    degree-field-label-en , figure-prefix-en , table-prefix-en , keywords-separator-en ,
    school-en , school-field-label-en , specialty-field-label-en , supervisor-field-label-en ,
    degree-level-en , degree-level-attribute , date-field-label-front ,
    graduate-affiliation-field-label , graduate-affiliation-major-field-label ,
    graduate-major-field-label , graduate-date-field-label ,
    graduate-enroll-date-field-label , graduate-student-id-field-label ,
    graduate-student-field-label , graduate-supervisor-field-label ,
    graduate-thesis-field-label , harbin-bachelor-school-bottom-mark ,
    harbin-school-bottom-mark , campus , heading-spread , hi , index-heading-en , index-heading-zh ,
    classified-index-international-field-label-zh , classified-index-national-field-label-zh ,
    classified-index-international-field-label-en , classified-index-national-field-label-en ,
    postdoc-udc-field-label , postdoc-author-field-label ,
    postdoc-classified-index-field-label , postdoc-report ,
    postdoc-research-end-date-field-label , postdoc-work-finish-date-field-label ,
    postdoc-specialty-field-label , postdoc-number-field-label ,
    postdoc-station-field-label , postdoc-secret-level-field-label , postdoc-separator ,
    postdoc-research-start-date-field-label , postdoc-report-submit-date-field-label ,
    postdoc-co-supervisor-field-label , publication-heading-zh , publication-heading-en ,
    resolution-heading-zh , resolution-heading-en , resume-heading-zh , resume-heading-en ,
    school-id , school-id-field-label , bachelor-second-authorization-text ,
    secret-level-field-label , shenzhen-doctor-midterm-note ,
    bachelor-affiliation-field-label-shenzhen ,
    bachelor-affiliation-field-label-weihai , bachelor-major-field-label ,
    bachelor-student-id-field-label , bachelor-supervisor-field-label ,
    bachelor-thesis-field-label , inline-separator , stage-document-type , stage-midterm ,
    stage-opening , supervisor-signature-field-label , title-separator-zh ,
    title-separator-en
  }
%    \end{macrocode}
%
% \changes{v3.2a}{0000/00/00}{参考文献的五个选项迁入 \texttt{hit/bib} 族，
%   改用连字符命名。旧名转回 \texttt{hit/class} 族处理，行为不变}
% \changes{v3.2a}{0000/00/00}{目录的四个选项迁入 \texttt{hit/toc} 族，改用连字符
%   命名，其中三个改为三态取值}
% \changes{v3.2a}{0000/00/00}{排版开关迁入 \texttt{hit/style} 族并改用连字符命名}
% 三个族的键。排版开关都是普通布尔，用 \cs{\_\_hit\_style\_pair:nn} 成对声明，
% 新名设的是旧标志位本身，所以两套名字天然一致。目录的前三个键用三态取值，
% \texttt{auto} 的含义分别是：\option{english} 博士排、其余不排；\option{blank}
% 英文版不插、其余插；\option{four} 一律不开。参考文献的五个键此前分散在
% \file{book} 与 \file{art} 两个模块、命名也不统一（\texttt{citetwo} 看不出是
% 控制什么的），新旧两套名字设的是同一批内部变量，所以可以混用。
%    \begin{macrocode}
\keys_define:nn { hit / style }
  {
    ziju .choices:nn = { regu , word }
      {
        \bool_gset_false:N \g__hit_ziju_regu_bool
        \bool_gset_false:N \g__hit_ziju_word_bool
        \bool_gset_true:c { g__hit_ziju_ \l_keys_choice_tl _bool }
      } ,
    ziju .value_required:n = true ,
    headrule .code:n = { \__hit_tristate_set:Nn \g__hit_head_rule_tl {#1} } ,
    pagenumberline .meta:nn = { hit / style / pagenumberline } {#1} ,
  }
\keys_define:nn { hit / style / pagenumberline }
  {
    mainmatter  .code:n = { \__hit_tristate_set:Nn \g__hit_page_number_line_main_tl  {#1} } ,
    frontmatter .code:n = { \__hit_tristate_set:Nn \g__hit_page_number_line_front_tl {#1} } ,
  }
\clist_map_inline:nn
  {
    { chapter-bold } { chapterbold } ,
    { chapter-hang } { chapterhang } ,
    { sans-title } { arialtitle } ,
    { caption-center-last } { capcenterlast } ,
    { subcaption-center-last } { subcapcenterlast } ,
    { abstract-upper } { absupper } ,
    { etitle-auto } { etitleauto } ,
    { glue } { glue } ,
    { open-right } { openright } ,
    { library } { library } ,
    { ragged-bottom } { raggedbottom } ,
  }
  { \__hit_style_pair:nn #1 }
%    \end{macrocode}
%
%    \begin{macrocode}
\keys_define:nn { hit / toc }
  {
    english .code:n = { \__hit_tristate_set:Nn \g__hit_toc_english_tl {#1} } ,
    blank   .code:n = { \__hit_tristate_set:Nn \g__hit_toc_blank_tl   {#1} } ,
    four    .code:n = { \__hit_tristate_set:Nn \g__hit_toc_four_tl    {#1} } ,
    sans-font .choice: ,
    sans-font / true  .code:n = { \bool_gset_true:N \g__hit_arialtoc_bool } ,
    sans-font / false .code:n = { \bool_gset_false:N \g__hit_arialtoc_bool } ,
    sans-font .default:n = true ,
  }
%    \end{macrocode}
%
%    \begin{macrocode}
\keys_define:nn { hit / bib }
  {
    adjacent-cite .choices:nn =
      { comma , endash }
      {
        \bool_gset_false:N \g__hit_citetwo_comma_bool
        \bool_gset_false:N \g__hit_citetwo_endash_bool
        \bool_gset_true:c { g__hit_citetwo_ \l_keys_choice_tl _bool }
      } ,
    adjacent-cite .value_required:n = true ,
    label-align .code:n = { \str_set:Nn \g_bib_labelalign_str {#1} } ,
    author-case .code:n = { \str_if_eq:nnT {#1} { full }
      { \bool_set_true:N \g_bib_authorupper_bool } } ,
    max-author .code:n =
      {
        \int_set:Nn \g_bib_maxauthor_int {#1}
        \int_if_zero:nT { \g_bib_maxauthor_int } { \int_gdecr:N \g_bib_maxauthor_int }
      } ,
    split-item .choice: ,
    split-item / true  .code:n = { \bool_gset_true:N \g__hit_splitbibitem_bool } ,
    split-item / false .code:n = { \bool_gset_false:N \g__hit_splitbibitem_bool } ,
    split-item .default:n = true ,
  }
%    \end{macrocode}
%
% \changes{v3.2a}{0000/00/00}{加旧名到新名的映射表与废弃警告。旧名照常生效，
%   只是编译时提示新写法}
% 兼容期的处理。旧名一律照常生效，只在日志里给一条提示。三类提示分开：
% 改过名的选项逐个报，元信息字段的扁平写法合并成一条报，
% \cs{documentclass} 里传排版类选项单独报。
%    \begin{macrocode}
\prop_gset_from_keyval:Nn \g__hit_renamed_option_prop
  {
    arialtoc         = { toc = { sans-font = ... } } ,
    tocfour          = { toc = { four = ... } } ,
    tocblank         = { toc = { blank = ... } } ,
    engtoc           = { toc = { english = ... } } ,
    arialtitle       = { style = { sans-title = ... } } ,
    chapterbold      = { style = { chapter-bold = ... } } ,
    chapterhang      = { style = { chapter-hang = ... } } ,
    capcenterlast    = { style = { caption-center-last = ... } } ,
    subcapcenterlast = { style = { subcaption-center-last = ... } } ,
    absupper         = { style = { abstract-upper = ... } } ,
    etitleauto       = { style = { etitle-auto = ... } } ,
    glue             = { style = { glue = ... } } ,
    openright        = { style = { open-right = ... } } ,
    library          = { style = { library = ... } } ,
    raggedbottom     = { style = { ragged-bottom = ... } } ,
    zijv             = { style = { ziju = ... } } ,
    bsheadrule       = { style = { headrule = ... } } ,
    bsmainpagenumberline  = { style = { pagenumberline = { mainmatter = ... } } } ,
    bsfrontpagenumberline = { style = { pagenumberline = { frontmatter = ... } } } ,
    citetwo          = { bib = { adjacent-cite = ... } } ,
    splitbibitem     = { bib = { split-item = ... } } ,
    biblabelalign    = { bib = { label-align = ... } } ,
    bibauthorupper   = { bib = { author-case = ... } } ,
    bibmaxauthor     = { bib = { max-author = ... } } ,
  }
%    \end{macrocode}
%
% \changes{v3.2a}{0000/00/00}{元信息字段的旧名走废弃期：仍然可用，用一次提示一次}
% 旧名到新名的对照。旧名在 v3.1e 就是公开接口，写在用户的 \file{cover.tex} 里，
% 所以保留到 v4.1。
%    \begin{macrocode}
\hit@define@legacy@terms
  {
    statesecrets=secret-level ,
    natclassifiedindex=classified-index-national ,
    intclassifiedindex=classified-index-international ,
    ctitle=title-zh ,
    cxueke=discipline-zh ,
    cauthor=author-zh ,
    csupervisor=supervisor-zh ,
    cassosupervisor=associate-supervisor-zh ,
    ccosupervisor=co-supervisor-zh ,
    caffil=affiliation-zh ,
    csubject=specialty-zh ,
    firstpagecdate=cover-date-zh ,
    cdate=date-zh ,
    cnumber=report-number ,
    cpositionname=station-zh ,
    cfinishdate=work-finish-date-zh ,
    csubmitdate=report-submit-date-zh ,
    cstartdate=research-start-date-zh ,
    cenddate=research-end-date-zh ,
    cstudentid=student-id ,
    cstudenttype=student-type-zh ,
    etitle=title-en ,
    exueke=discipline-en ,
    eauthor=author-en ,
    esupervisor=supervisor-en ,
    eassosupervisor=associate-supervisor-en ,
    ecosupervisor=co-supervisor-en ,
    eaffil=affiliation-en ,
    esubject=specialty-en ,
    edate=date-en ,
    estudenttype=student-type-en ,
    ckeywords=keywords-zh ,
    ekeywords=keywords-en
  }
\hit@define@legacy@title{ctitlecover}{zh}{first}
\hit@define@legacy@title{ctitleone}{zh}{first}
\hit@define@legacy@title{ctitletwo}{zh}{second}
\hit@define@legacy@title{csubtitle}{zh}{subtitle}
\hit@define@legacy@title{esubtitle}{en}{subtitle}
\prop_gset_from_keyval:Nn \g__hit_renamed_load_option_prop
  {
    font     = { font-en = ... } ,
    cjk-font = { font-zh = ... } ,
  }
%    \end{macrocode}
%
% \changes{v3.2a}{0000/00/00}{提示文案、只报一次的机制与类选项扫描挪到
%   \file{src/utils/hit-meta.dtx}}
% 单双面由学位与校区决定，所以扫描类选项时额外查一下 \option{oneside} 与
% \option{twoside}。另外两个类没有这条。
%    \begin{macrocode}
\cs_set_protected:Npn \__hit_class_option_extra:
  {
    \clist_if_in:nVT { oneside , twoside } \l__hit_warn_message_tl
      {
        \exp_args:NV \__hit_warn_once:nnn \l__hit_warn_message_tl
          { manual-sides } { }
      }
  }
\__hit_check_class_options:
\keys_define:nn { hit }
  {
    info  .meta:nn = { hit / info  } {#1} ,
    style .meta:nn = { hit / style } {#1} ,
    toc   .meta:nn = { hit / toc   } {#1} ,
    bib   .meta:nn = { hit / bib   } {#1} ,
    const .meta:nn = { hit / const } {#1} ,
    structure .meta:nn = { hit / structure } {#1} ,
    theorem .code:n = { \hit@theorem@style {#1} } ,
    theorem-list .code:n = { \hit@theorem@list {#1} } ,
    theorem-star .code:n = { \hit@theorem@declare@star@list {#1} } ,
    unknown .code:n = { \__hit_legacy_key:n {#1} } ,
  }
\RenewDocumentCommand \hitsetup { +m } { \keys_set:nn { hit } {#1} }
%    \end{macrocode}
%
% \changes{v3.1d}{2025/03/03}{将 \cs{firstpagecdate} 的设置藏起来}
%    \begin{macrocode}
\AfterPreamble{%
  \ifdefempty{\hit@cover@date@zh}{\use:c{cover-date-zh}{\hit@date@zh}}{}%
}
%</thesis-meta>
%    \end{macrocode}
%
% \Finale
